\documentclass[12pt]{article}
\usepackage[utf8]{inputenc}
\usepackage[T1]{fontenc}
\usepackage{hyperref}
\title{Aribwaung Language Resources for K-12 Curriculum Development}
\author{Compiled by Claude Sonnet 5.5 (claude-sonnet-5-5) for Andr\'as Kornai}
\date{}
\begin{document}
\maketitle

\section{Introduction}
Aribwaung (also called Yalu or Jaloc; ISO 639-3 ylu; Glottocode arib1240) is an Austronesian language of the Markham group, Huon Gulf branch of Oceanic. It is spoken in the village of Yalu, Wampar Rural LLG, Morobe Province, Papua New Guinea. Wikipedia lists about 500 speakers (2008), while Global Recordings Network gives a population of 1,000. It should not be confused with Aribwatsa (Lae), which is a different and extinct language.

\section{Primary Educational Resources}
I found no primers, textbooks or school programme for Aribwaung.

\section{Linguistic and Reference Materials}
\begin{itemize}
\item The English Wikipedia article gives a general overview. \url{https://en.wikipedia.org/wiki/Aribwaung_language}
\item The ANU Archives hold an item described as ``Aribwaungg (Yalu) language Papua New Guinea'', which includes a lexis, a syntactic check list and notes. It also refers to Malcolm Ross's grammar sketch and text. I saw only the title. \url{https://archivescollection.anu.edu.au/index.php/aribwaungg-yalu-language-papua-new-guinea-includes-lexis-syntactic-check-list-notes-yalu-from-malcolm-ross-grammar-sketch-and-text}
\item The Max Planck EVA numeral-systems page for Aribwaung (Yalu) records its counting system. The system is base-two with hand and foot tallying, and the word \emph{artsamo} means 20. \url{https://lingweb.eva.mpg.de/channumerals/Aribwaung-Yalu.htm}
\end{itemize}

\section{Religious and Community Resources}
Global Recordings Network has a language page for Aribwaung. I saw only the listing and did not check what recordings it offers. \url{https://globalrecordings.net/en/language/ylu}

\section{Digital and Audio Resources}
PARADISEC catalogues Aribwaung items, recorded in Morobe Province during Susanne Holzknecht's Markham fieldwork in 1985--88. I did not check access conditions. The catalogue entries include:
\begin{itemize}
\item \url{https://catalog.paradisec.org.au/repository/SH2/ANUA520_4401}
\item \url{https://catalog.paradisec.org.au/repository/SH2/ANUA520_4407}
\item \url{https://catalog.paradisec.org.au/repository/SH2/ANUA520_4109}
\item \url{https://catalog.nabu-prod.paradisec.org.au/collections/SH2/items/ANUA520_4409}
\end{itemize}

\section{Curriculum Development Context}
The language is still spoken, but I found no school teaching materials. The archival recordings and the lexis and grammar-sketch materials could support community-led development of materials. A curriculum would need to be built with the Yalu community.
\end{document}